% Generated draft: requires mathematical and source review.
\documentclass[11pt]{article}
\usepackage[margin=25mm]{geometry}
\usepackage{fontspec}
\setmainfont{Latin Modern Roman}
\usepackage{amsmath,amssymb,mathtools,unicode-math}
\setmathfont{Latin Modern Math}
\usepackage{enumitem,needspace,xurl,hyperref}
\setlength{\parindent}{0pt}
\setlength{\parskip}{5pt}
\setlength{\emergencystretch}{3em}
\newcommand{\N}{\mathbb{N}}
\newcommand{\Z}{\mathbb{Z}}
\newcommand{\Q}{\mathbb{Q}}
\newcommand{\R}{\mathbb{R}}
\newcommand{\C}{\mathbb{C}}
\newcommand{\F}{\mathbb{F}}
\newcommand{\A}{\mathbb{A}}
\newcommand{\PP}{\mathbb{P}}
\newcommand{\E}{\mathbb{E}}
\newcommand{\eps}{\varepsilon}
\newcommand{\dd}{\,\mathrm{d}}
\begin{document}
\section*{UnsolvedMath ID: 20002700}
\textbf{Source problem code:} catalogue entry
\subsection{Definitions and mathematical statement}
The source asks to explain cluster algebras of infinite type, but ends with an incomplete ``See'' and gives no reference or precise classification target. Infinite cluster type at finite rank, an infinite mutation class of exchange matrices, and infinite rank are different notions. The aid supplies a complete comparison only in coefficient-free rank two.

Fix positive integers $b,c$ and initial principal exchange matrix
\[
 B=\begin{pmatrix}0&b\\-c&0\end{pmatrix}.
\]
There are no frozen variables here. Let $x_1,x_2$ be algebraically independent and generate the rank-two cluster algebra inside $\Q(x_1,x_2)$ by the bi-infinite recurrence
\[
 x_{m-1}x_{m+1}=
 \begin{cases}1+x_m^b,&m\text{ odd},\\
               1+x_m^c,&m\text{ even}.
 \end{cases}
\]
This convention corresponds to mutating the first initial variable first. The source aid reports that the labeled principal-matrix mutation class is always exactly $\{B,-B\}$, while the number of cluster variables is finite exactly for $bc\le3$. Thus two exchange matrices do not imply finite cluster type.

For $bc\ge4$, write the reduced Laurent expression of a noninitial variable as
\[
 x_j=\frac{P_j(x_1,x_2)}{x_1^{d_{j,1}}x_2^{d_{j,2}}},
\]
with numerator not divisible by either initial variable. For the forward alternating-mutation ray the exact reported denominator vectors are
\[
 d_{2k+3}=C^k\binom10,\qquad
 d_{2k+4}=C^k\binom b1,\qquad k\ge0,
 \quad C=\begin{pmatrix}bc-1&-b\\c&-1\end{pmatrix}.
\]
The opposite ray has the corresponding exchanged-label description. The index convention matters; these are not asserted to list every initial or reverse-ray variable with the same positive index formula.

At the affine threshold $bc=4$, $(C-I)^2=0$ and $C^k=I+k(C-I)$, so both vector sequences grow linearly in $k$. For $bc>4$ they grow exponentially with per-two-mutation spectral radius
\[
 \Lambda=\frac{bc-2+\sqrt{bc(bc-4)}}2>1.
\]
In both infinite regimes the exchange graph is a bi-infinite path, so every radius-$R$ graph ball has exactly $2R+1$ vertices for integer $R\ge0$. Consequently exchange-graph ball growth and principal mutation-class size do not distinguish the affine and indefinite rank-two cases, whereas denominator growth does. This comparison does not supply a theory or classification for all infinite-type or infinite-rank cluster algebras.
\subsection{Short English statement}
Clarify the different meanings of infinite type. In coefficient-free rank two, affine and indefinite cases both have two principal exchange matrices and a path exchange graph, but their denominator vectors grow linearly and exponentially respectively.
\subsection{Sources}
\begin{itemize}
\item American Institute of Mathematics workshop \emph{Braid groups, clusters and free probability}, Problem 5.1, \url{https://aimath.org/WWN/braidgroups/braidgroups.pdf}.
\item Local UnsolvedMath registry, permanent ID 20002700, source code \texttt{catalogue entry}, canonical location \texttt{aim-probability-notes.json}, \texttt{notes[141]}, record dated 14 August 2026. The original statement ends with ``See'' without a citation; the original literature field is empty.
\item Supplied research attempt 1: partial result, problem status at run partially solved, rigor labelled proved, novelty confidence low. Its rank-two three-invariant comparison is machine-generated research aid, not peer reviewed; independently verify before citation. No separate authoritative reference for the classification or formulas is supplied.
\end{itemize}
\end{document}
